\chapter{Conclusiones}\label{cap:con}

El trabajo partía de una pregunta: si se puede reducir de forma apreciable el coste de ejecutar un agente DRL+GNN para ingeniería de tráfico sin empeorar el enrutamiento que produce. Para responderla se puso en marcha el agente de ENERO con bibliotecas actuales, se entrenó un agente base con una recompensa corregida, se le aplicaron poda, cuantización y destilación, y los modelos resultantes se evaluaron en tres topologías que no habían intervenido en ningún momento del proceso y, después, en siete redes mayores. Para valorar el alcance práctico de los resultados, ENERO, original y comprimido, se comparó además con la formulación exacta del mismo problema. Este capítulo recoge lo que se ha aprendido (Sección~\ref{sec:con_conclusiones}), con qué límites (Sección~\ref{sec:con_limitaciones}), qué quedaría por hacer (Sección~\ref{sec:con_futuro}) y una valoración personal (Sección~\ref{sec:con_valoracion}).

% =====================================================================
\section{Conclusiones}\label{sec:con_conclusiones}

La respuesta es que sí: el coste del agente se puede reducir de forma importante sin empeorar el enrutamiento, aunque cada técnica lo consigue en un aspecto distinto y su beneficio depende del hardware y del tipo de red. Sin embargo, la comparación con la formulación exacta lleva a una conclusión de mayor alcance, que se recoge al final de esta sección: en las redes estudiadas, ni ENERO ni sus versiones comprimidas son la mejor forma de resolver este problema.

La primera conclusión es que \textbf{la compresión conserva la calidad del enrutamiento}. En las 150 matrices de las tres topologías de evaluación, todos los modelos comprimidos, seguidos de la LS, mejoran la MLU de OSPF entre un 32,6 y un 35,0\,\% de media, frente al 34,3\,\% del modelo original, y todos superan a la LS aplicada sin agente. El modelo podado al 50\,\% y las variantes FP16 e INT8 de rango dinámico no presentan diferencias estadísticamente significativas con el original, y el estudiante obtiene en media la misma calidad; solo la poda extrema, al 90\,\%, es ligeramente peor, con diferencias estadísticamente significativas. El resultado se mantiene en las siete redes mayores, donde los modelos comprimidos siguen rindiendo como el original.

La segunda es que \textbf{la destilación es la técnica más adecuada para un procesador convencional}. El estudiante tiene 4,6 veces menos parámetros y es el único modelo que acelera la inferencia, 3,2 veces, porque reduce la arquitectura y no solo la forma de guardar los pesos. Conserva en media la calidad del original y reduce entre un 5 y un 25\,\% el tiempo y la energía de cada optimización, según la calidad del punto de partida que entregue a la LS. Al repetir la destilación con varias semillas se vio que el resultado de un estudiante concreto varía de un entrenamiento a otro, por lo que conviene validarlo en la red de destino antes de desplegarlo. El estudio de configuraciones mostró también que, para destilar este agente, conviene comparar solo las distribuciones de probabilidad del maestro y del estudiante, con una temperatura alta, y no sus valores numéricos.

La tercera es que \textbf{la poda y la cuantización reducen el tamaño, pero no el tiempo, en un procesador convencional}. El agente tiene bastante redundancia: casi un tercio de sus pesos se pueden eliminar sin reentrenar y, reentrenando, el modelo podado a la mitad conserva el 98\,\% de la mejora del original y ocupa un 38\,\% menos comprimido. Pero la poda no estructurada solo acelera con hardware que aproveche los ceros. La cuantización con TensorFlow Lite reduce el fichero entre un 25 y un 40\,\%, aunque en un modelo tan pequeño el grafo de operaciones ocupa más que los pesos, y en este procesador resulta más lenta que TensorFlow porque el modelo convertido se ejecuta operación a operación.

La cuarta es que \textbf{en ENERO la red neuronal no es el cuello de botella}. La LS ocupa en torno al 70\,\% del tiempo de optimización, y dentro de la etapa DRL pesa más construir los grafos de los candidatos y simular la red que la inferencia. Comprimir el agente mejora sobre todo la parte que depende de él, y la calidad de la solución que entrega a la LS influye tanto como su velocidad.

La quinta conclusión es que \textbf{el agente aporta más en redes parecidas a las de entrenamiento}. En las topologías con una estructura similar, ENERO mejora claramente a la LS sin agente y queda cerca del óptimo de su propio espacio de búsqueda; en redes mayores y más dispersas, su aportación es menor y se aleja mucho de ese óptimo. Si se utiliza ENERO, la recomendación práctica es entrenar o ajustar el agente con topologías del tipo en que se va a desplegar y comprimirlo después: con poda moderada si importa el almacenamiento, o con destilación si importa la latencia.

La sexta es que \textbf{el acuerdo de decisiones no sirve para medir la calidad de un agente comprimido}. Modelos que coinciden con el original en una de cada cinco decisiones obtienen una calidad similar, porque hay muchas formas casi igual de buenas de desviar el tráfico. A diferencia de lo habitual al comprimir clasificadores, la evaluación debe hacerse sobre la tarea completa.

Por último, la conclusión de mayor alcance se refiere a la metodología. Tras evaluar ENERO y sus versiones comprimidas no solo por la calidad del enrutamiento, sino también por su aplicabilidad práctica (el tiempo de cada optimización, el coste de entrenar y ajustar el agente y su capacidad de generalizar a redes distintas), \textbf{la metodología más adecuada para redes de hasta unas decenas de nodos es la formulación exacta del enrutamiento por punto intermedio como MILP}. Resuelto sobre las mismas demandas críticas que el agente, el MILP obtiene el óptimo del problema que aborda ENERO en menos de un segundo de media, unas 70 veces más rápido, con mejor resultado que ENERO original en las diez redes evaluadas y sin entrenamiento previo. Resuelto sobre todas las demandas, alcanza prácticamente la cota de programación lineal en unos segundos. El aprendizaje por refuerzo profundo, en cambio, requiere días de entrenamiento, generaliza mal a redes distintas de las de entrenamiento y, precisamente en esas redes, queda lejos del óptimo. La compresión reduce el coste del agente, pero no cambia esta conclusión, porque la desventaja de ENERO no está en el tamaño de su red neuronal, sino en el enfoque. Su posible interés queda limitado a redes mucho mayores, en las que el método exacto deje de resolverse en un tiempo razonable, algo que este trabajo no ha podido comprobar.

En cuanto a los objetivos de la Sección~\ref{sec:int_objetivos}, se ha revisado el estado del arte, se ha migrado ENERO a Gymnasium y entrenado un agente base en CPU, se ha corregido la recompensa, se han aplicado las tres familias de compresión y se han evaluado con una metodología que separa el coste de la red del de la optimización completa e incluye la cota de programación lineal, el acuerdo de decisiones y contrastes estadísticos emparejados. También se ha estimado el impacto energético de cada alternativa y se ha comparado ENERO, original y comprimido, con la formulación exacta del problema, en calidad y en aplicabilidad práctica.

% =====================================================================
\section{Limitaciones}\label{sec:con_limitaciones}

Las conclusiones anteriores deben leerse teniendo en cuenta las siguientes limitaciones:

\begin{itemize}
	\item \textbf{Simulación y hardware}. Todo se ha hecho en simulación y en un único procesador de propósito general. En hardware con soporte para operaciones dispersas o aritmética entera, la poda y la cuantización podrían comportarse de otra forma, y la cuantización solo se ha probado con TensorFlow Lite.
	\item \textbf{Pocas topologías de entrenamiento}. El agente se entrena con tres topologías, lo que limita su generalización a redes de otro tipo. Las siete redes mayores se evaluaron con 10 matrices cada una, de modo que las conclusiones sobre ellas son orientativas. Esta limitación se debe al tiempo disponible.
	\item \textbf{Variabilidad entre entrenamientos}. Solo se ha medido para la destilación, con cuatro semillas, porque cada repetición ocupa el equipo varias horas. El agente base y la poda corresponden a un único entrenamiento.
	\item \textbf{Ajuste limitado de hiperparámetros}. Los hiperparámetros se fijaron tras pocas pruebas por el coste de cada entrenamiento en CPU.
	\item \textbf{Tamaño de las redes y otros optimizadores}. Las redes evaluadas tienen como máximo 43 nodos. En ellas el MILP se resuelve en segundos, pero no se ha comprobado cómo crece su tiempo en redes de cientos de nodos, que es donde un agente entrenado podría compensar. Tampoco se ha comparado con otros optimizadores, como DEFO, y el tiempo medido del MILP no incluye la construcción del modelo.
	\item \textbf{Estimación energética}. La energía se estima con la potencia media de la CPU bajo carga y los tiempos medidos, sin medir el consumo durante cada evaluación.
\end{itemize}

% =====================================================================
\section{Trabajo futuro}\label{sec:con_futuro}

De los resultados y de las limitaciones se desprenden varias líneas de trabajo:

\begin{itemize}
	\item \textbf{Entrenar para la red de destino}. Entrenar o ajustar el agente con topologías del tipo en que se va a desplegar, y comprobar que la compresión conserva esa mejora, es el paso más directo para aumentar su aportación en redes grandes o dispersas.
	\item \textbf{Combinar técnicas}. Podar o cuantizar el estudiante destilado podría acumular beneficios, y destilar hacia estudiantes aún más pequeños permitiría localizar el límite de capacidad del modelo.
	\item \textbf{Poda estructurada}. Eliminar neuronas o dimensiones completas del estado, en lugar de pesos sueltos, produce modelos densos más pequeños que sí se aceleran en cualquier hardware.
	\item \textbf{Otros motores de inferencia}. Probar la cuantización con motores que fusionen las operaciones del grafo, o en hardware con aritmética entera, donde debería aparecer su ventaja en velocidad.
	\item \textbf{Más topologías y semillas}. Repetir con varias semillas el agente base y la poda, y buscar un criterio de validación que permita anticipar qué estudiante generalizará mejor.
	\item \textbf{Escalabilidad del MILP y enfoques híbridos}. Medir el tiempo del MILP en redes de cientos de nodos, para localizar el tamaño a partir del cual un agente entrenado podría compensar, y estudiar enfoques que combinen ambos métodos: usar la solución del agente como punto de partida del MILP, o usar las soluciones del MILP para entrenar o destilar el agente.
	\item \textbf{Acelerar el resto del optimizador}. Como la red no es el cuello de botella, hay margen en la construcción de los grafos, la simulación de la red y la LS, por ejemplo vectorizándolas o limitando la búsqueda a los candidatos mejor valorados por el agente.
	\item \textbf{Porcentaje de demandas críticas}. Estudiar cómo influye el porcentaje de demandas críticas en la calidad y en el tiempo de optimización según el tamaño de la red, y si conviene adaptarlo a cada topología.
\end{itemize}

% =====================================================================
\section{Valoración personal}\label{sec:con_valoracion}

Cuando empecé el trabajo apenas conocía el aprendizaje por refuerzo y las redes neuronales de grafos (solo por encima), y no tenía claro cómo encajaban en la gestión de una red. Entender ENERO desde el artículo hasta el código me obligó a juntar cosas que en el grado había visto en asignaturas distintas: enrutamiento, optimización y aprendizaje automático.

Lo que más tiempo me ha llevado no ha sido la compresión, sino conseguir que todo lo demás funcionara bien. Algunos errores no daban ningún fallo visible: la primera evaluación de los modelos de TensorFlow Lite les pasaba un solo grafo en lugar de los de todos los candidatos, y el agente elegía siempre la primera acción sin que nada lo indicara; el sesgo de la recompensa solo apareció al revisar el entorno línea a línea. Corregirlos obligó a repetir buena parte del trabajo en septiembre y octubre, pero ahora puedo defender cada cifra de la memoria.

También he tenido que aprender a planificar con experimentos que duran días. Con un portátil con GPU no compatible, con el procesador a unos 95~$^\circ$C durante días enteros, cada entrenamiento había que pensarlo bien antes de lanzarlo, trocearlo para poder reanudarlo cuando se cortaba y repartir el trabajo entre dos equipos.

Por último, me quedo con que algunos resultados no salieron como esperaba. El primer estudiante destilado parecía mejor que su maestro, y solo al repetir la destilación con otras semillas vi que el resultado variaba de un entrenamiento a otro; me ha enseñado a no fiarme de un único entrenamiento. Tampoco esperaba que la red neuronal no fuera el cuello de botella del optimizador, y eso me hizo entender que, para valorar una técnica de compresión, hay que medirla dentro del sistema completo y no solo sobre el modelo.

Lo que más me ha hecho reflexionar ha sido la comparación con el MILP. Empecé el trabajo dando por hecho que un agente entrenado era la forma más moderna y eficaz de resolver este problema, y me costó aceptar un resultado que no favorecía al método que estaba estudiando. Me ha enseñado que una comparación entre métodos tiene que hacerse sin ideas preconcebidas, y que un resultado negativo, si está bien medido, tiene tanto valor como uno positivo.
